\documentclass[11pt]{article}

\usepackage{amsmath,amssymb,amsthm}
\usepackage[margin=1.1in]{geometry}
\usepackage{bm}
\usepackage[colorlinks=true,linkcolor=blue,citecolor=blue]{hyperref}

\newtheorem{proposition}{Proposition}
\newtheorem{remark}{Remark}
\newcommand{\R}{\mathbb{R}}
\newcommand{\eps}{\varepsilon}
\newcommand{\ip}[2]{\langle #1, #2\rangle}
\newcommand{\norm}[1]{\lVert #1 \rVert}

\title{Gauge fixing by penalty and iterated refinement,\\
       and the sliding-interface displacement discontinuity:\\
       the regularisation machinery of the gauged formulation}
\author{mfemElasticity method notes}
\date{}

\begin{document}
\maketitle

\noindent
Companion to \texttt{doc/gauged\_fluid.md} and
\texttt{doc/gravitating\_elasticity.md}~\S3.1, which state \emph{what} is
gauged; this note explains \emph{how} the gauge freedom is handled
numerically. The technique is a small-parameter penalty combined with
what the inverse-problems literature calls \emph{iterated Tikhonov
regularisation}; nothing here is original, but the pieces are scattered,
and the failure mode we met with the vacuum extension is instructive
enough to deserve a self-contained account.

\section{The problem: a gauged discrete system}

After discretisation (and after projecting out the rigid-body null
space, which is handled separately and exactly by the null-space
projectors), the linear problems of interest have the form
\begin{equation}
  A\,x = b, \qquad A = A^{T} \succeq 0,
  \label{eq:system}
\end{equation}
where $A \in \R^{n\times n}$ is the assembled physical operator and $b$
is compatible, $b \perp \ker A$. The distinguishing feature is a large
\emph{gauge subspace}
\[
  \mathcal{G} \;\subseteq\; \ker A
  \quad\text{(or nearly so),}
\]
of dimension comparable to $n$ itself, along which the physics is
indifferent: any two solutions differing by an element of $\mathcal{G}$
represent the same physical state, and every observable is a
gauge-invariant functional. Two instances in the code:

\begin{itemize}
\item \textbf{The gauged fluid.} In the relabelled fluid formulation the
  fluid regions keep a displacement field, but only its volumetric
  content (through $\kappa\,(\nabla\!\cdot u)$) and the gravity coupling
  are physical; the deviatoric part of the fluid motion is pure
  relabelling. $\mathcal{G}$ is (essentially) the space of
  divergence-free fluid-supported fields --- huge, mesh-dependent, and
  not aligned with any set of nodal degrees of freedom, so it cannot be
  eliminated by constraining dofs.
\item \textbf{The ball-wide vacuum extension.} In the referential
  gravity formulation the linearised perturbation needs a displacement
  extension through the buffer shell; the extension is entirely gauge
  (any smooth field with the correct trace on $\partial B$ gives the
  same observables), and if the extension is carried as unknowns, the
  buffer block of $A$ has no stiffness at all.
\end{itemize}

Solving \eqref{eq:system} directly is unattractive: Krylov methods on a
semi-definite operator with an $O(n)$-dimensional kernel converge to
\emph{a} solution only if the preconditioner respects the kernel, and
standard AMG/Gauss--Seidel preconditioners do not. The penalty route
restores definiteness first and removes its bias afterwards.

\section{Penalty regularisation}

Choose a symmetric positive semi-definite \emph{penalty operator}
$Q = Q^{T} \succeq 0$ that is positive definite on the gauge subspace,
\begin{equation}
  x \in \mathcal{G}, \; x \neq 0 \quad\Longrightarrow\quad x^{T} Q x > 0 ,
\end{equation}
and replace \eqref{eq:system} by
\begin{equation}
  A_{\eps}\, x_{\eps} = b, \qquad A_{\eps} := A + \eps\, Q ,
  \label{eq:penalised}
\end{equation}
with a small parameter $\eps > 0$. In the code $Q$ is assembled from an
ordinary bilinear form:
\begin{itemize}
\item \emph{Deviatoric penalty} (gauged fluid):
  $q(u,v) = \int 2\mu_{g}\, \mathrm{dev}\,\varepsilon(u) :
  \mathrm{dev}\,\varepsilon(v)$ on the fluid attributes. This is the
  natural choice because it \emph{vanishes identically on the physical
  (volumetric) fluid response}: the penalty is invisible to the physics
  it accompanies, up to the coupling terms.
\item \emph{Harmonic penalty} (extension fields):
  $q(u,v) = \int \mu_{g}\, \nabla u : \nabla v$ on the buffer
  attributes, which penalises everything --- there is no physical motion
  in a vacuum buffer to spare.
\end{itemize}
The reference scale $\mu_{g}$ (a bulk or shear modulus of the adjacent
solid) makes $\eps$ dimensionless.

\subsection{What one penalised solve gives}
\label{sec:onesolve}

Assume for the analysis that $Q \succ 0$ on the whole space (in
practice: on the complement of the projected rigid modes; on the
orthogonal complement of $\mathcal{G}$ any spd completion works and
none of the conclusions change). Introduce the generalised eigenpairs
\begin{equation}
  A\, v_{i} = \lambda_{i}\, Q\, v_{i}, \qquad
  \lambda_{1} \le \lambda_{2} \le \dots , \qquad
  v_{i}^{T} Q v_{j} = \delta_{ij} ,
  \label{eq:gevp}
\end{equation}
which simultaneously diagonalise $A$ and $Q$. Gauge directions are the
$\lambda_{i} = 0$ modes; the \emph{physical spectrum} is
$\{\lambda_{i} > 0\}$, and its smallest member
\begin{equation}
  \lambda_{\min} := \min \{\, \lambda_{i} : \lambda_{i} > 0 \,\}
  \label{eq:lmin}
\end{equation}
is the quantity the whole method hinges on. Expanding
$x = \sum_i y_i v_i$, $b_i := v_i^T b$ (compatibility: $b_{i} = 0$ on
gauge modes), the exact and penalised solutions are, mode by mode,
\begin{equation}
  y_{i} = \frac{b_{i}}{\lambda_{i}} \quad (\lambda_{i} > 0), \qquad
  y_{\eps,i} = \frac{b_{i}}{\lambda_{i} + \eps} .
\end{equation}
Hence:
\begin{proposition}
  The penalised solution $x_{\eps}$ (i) fixes the gauge exactly: its
  component along every gauge mode is zero, i.e.\ $x_\eps$ is the
  representative of the physical solution that minimises the penalty
  norm $x^{T} Q x$; and (ii) carries a relative bias on each physical
  mode of
  \begin{equation}
    \frac{y_{i} - y_{\eps,i}}{y_{i}}
      = \frac{\eps}{\lambda_{i} + \eps}
      \;\le\; \frac{\eps}{\lambda_{\min} + \eps} .
    \label{eq:bias}
  \end{equation}
\end{proposition}
So a single solve of \eqref{eq:penalised} produces a perfectly
gauge-fixed answer whose \emph{physics} is wrong by
$O(\eps/\lambda_{\min})$. The classical dilemma follows immediately:
shrinking $\eps$ shrinks the bias but sends the condition number
$\kappa(A_{\eps}) \sim \lambda_{\max}^{A}/(\eps\,\lambda_{\min}^{Q})$ to
infinity --- the preconditioner degrades exactly as the answer improves.
This is Tikhonov regularisation in its ridge-regression form, and the
dilemma is the textbook one.

\section{Iterated refinement}

The escape is to keep $\eps$ \emph{moderate} and remove the bias by
iteration instead. Define
\begin{equation}
  x^{(1)} = A_{\eps}^{-1} b, \qquad
  x^{(k+1)} = x^{(k)} + A_{\eps}^{-1}\bigl( b - A\, x^{(k)} \bigr),
  \qquad k = 1, 2, \dots
  \label{eq:iteration}
\end{equation}
Each step is one additional solve \emph{with the same operator}
$A_{\eps}$ --- the factorisation or AMG hierarchy and the Krylov setup
are built once and reused, so refinements are nearly free compared with
the first solve. The residual is computed with the \emph{physical}
operator $A$, which is the whole point: the iteration converges (where
it converges) to a solution of \eqref{eq:system}, not of
\eqref{eq:penalised}.

\subsection{Convergence, mode by mode}

Subtracting \eqref{eq:system} from \eqref{eq:iteration}, the error
$e^{(k)} = x^{(k)} - x$ obeys
\begin{equation}
  e^{(k+1)} = \bigl( I - A_{\eps}^{-1} A \bigr)\, e^{(k)}
            = \eps\, A_{\eps}^{-1} Q\, e^{(k)} ,
  \label{eq:errrec}
\end{equation}
and in the eigenbasis \eqref{eq:gevp} the iteration matrix is diagonal:
\begin{equation}
  e^{(k+1)}_{i} = \frac{\eps}{\lambda_{i} + \eps}\; e^{(k)}_{i}
  \qquad\Longrightarrow\qquad
  \frac{y_{i} - y^{(k)}_{i}}{y_{i}}
    = \Bigl( \frac{\eps}{\lambda_{i} + \eps} \Bigr)^{\!k} .
  \label{eq:contraction}
\end{equation}
Three consequences:
\begin{enumerate}
\item On \textbf{physical modes} the relative bias after $k$ solves is
  $\le (\eps/(\lambda_{\min}+\eps))^{k}$: with
  $\lambda_{\min} = O(1)$ and $\eps = 10^{-2}$, two or three
  refinements reach $10^{-4}$--$10^{-6}$, which is below discretisation
  error. The contraction factor per step is
  \begin{equation}
    \rho \;=\; \frac{\eps}{\lambda_{\min} + \eps}
    \;\approx\; \frac{\eps}{\lambda_{\min}} .
    \label{eq:rho}
  \end{equation}
\item On \textbf{gauge modes} ($\lambda_{i} = 0$, $b_{i} = 0$) the
  factor is $1$: whatever gauge value the first solve produced (namely
  zero, the $Q$-minimal representative) is preserved unchanged. The
  iteration never un-fixes the gauge.
\item The fixed point on the physical subspace is the exact solution of
  \eqref{eq:system}; the method is \emph{consistent} for every
  $\eps > 0$.
\end{enumerate}

In the inverse-problems literature \eqref{eq:iteration} is
\emph{(stationary) iterated Tikhonov regularisation}
\cite{Lardy1975,EnglHankeNeubauer}: $k$-times-iterated Tikhonov has
qualification $k$, i.e.\ bias $O(\eps^{k})$ instead of the single-step
$O(\eps)$. Numerical analysts will equally recognise it as \emph{defect
correction} (iterative refinement with a deliberately perturbed
operator): $A_\eps$ plays the role of the ``approximate inverse'' and
$A$ supplies the true residual. Both views are exact descriptions of
the same three lines of code.

\subsection{When it must fail: the spectral condition}
\label{sec:failure}

Everything above assumed \eqref{eq:lmin}: the smallest \emph{physical}
generalised eigenvalue of $(A, Q)$ stays bounded away from zero,
\emph{uniformly in the discretisation}. In words:
\begin{quote}
  \emph{On every physical mode, the penalty must be dominated by the
  physical operator, with a mesh- and order-independent margin.}
\end{quote}
This is a structural property of the pair $(A, Q)$, not a tuning
question, and the two applications in the code sit on opposite sides of
it.

\paragraph{Gauged fluid: the condition holds.}
Here $Q$ is a deviatoric elasticity form with modulus $\mu_{g}$ and $A$
contains genuine elastic stiffness (solid regions, and the fluid's own
$\kappa$-volumetric term) of the same differential order. Both scale
identically under mesh refinement, so the generalised spectrum is
$h$-independent and $\lambda_{\min} \sim \mu_{\mathrm{solid}}/\mu_{g}$
up to geometry-dependent constants. The observed contraction per
refinement is $\rho \approx \eps\,\mu_{g}/\mu_{\mathrm{solid}}$,
uniformly across meshes and orders --- the benign case.

\paragraph{Ball-wide vacuum extension: the condition fails.}
In the buffer the physical operator has \emph{no stiffness at all}:
$A$ restricted there consists of the zeroth-order (undifferentiated)
gravity couplings, which are $L^{2}$-bounded, while $Q$ is a harmonic
(first-derivative) form whose Rayleigh quotients on fine-scale buffer
modes grow like $h^{-2}$ (and with polynomial order). The generalised
eigenvalues of buffer-dominated modes therefore behave like
$\lambda \sim h^{2}$: $\lambda_{\min} \to 0$ under refinement, the
contraction factor \eqref{eq:rho} tends to $1$ and the amplification of
numerical error (below) blows up. This is exactly what was observed:
the biased single-solve regime works (with $O(\eps)$ error), but
refinement \emph{diverges} at order 2. No choice of $\eps$ or $k$
repairs it, because the failure is in the spectrum of the pair, not in
the parameters. The cure was structural: replace the buffer unknowns by
a \emph{prescribed} extension operator (option (b) of
\texttt{gravitating\_elasticity.md}~\S3.1), which removes the gauge
unknowns instead of penalising them. The penalty machinery returns for
the nonlinear problem, where the buffer motion re-acquires physical
meaning.

\subsection{Semi-convergence and the operating point}
\label{sec:semiconv}

Formula \eqref{eq:contraction} is exact arithmetic. In practice each
application of $A_{\eps}^{-1}$ is a Krylov solve to tolerance $\tau$,
and the residual $b - A x^{(k)}$ is formed in floating point. A
perturbation $\delta$ of the residual enters the next iterate as
$A_{\eps}^{-1}\delta$, and on modes with small $\lambda_{i} + \eps$
this is an amplification by up to $1/\eps$ (in the $Q$-normalised
variables). The bias therefore decays like $\rho^{k}$ only down to a
noise floor of order $\tau/\eps$, after which further refinements
random-walk or grow: the error as a function of $k$ is U-shaped. This
is the \emph{semi-convergence} familiar from iterative regularisation,
and it is why the recommended operating point is
\begin{equation}
  \eps \sim 10^{-2}, \qquad k = 2\text{--}3 \ \text{refinements},
  \qquad \tau \lesssim 10^{-10} ,
\end{equation}
rather than ``$\eps$ as small as possible''. Two further practical
notes:
\begin{itemize}
\item Modes adjacent to the projected null space (the near-rigid,
  Slichter-soft degree-1 responses) have small $\lambda_{i}$ of
  \emph{physical} origin; they converge at the slow rate
  $\eps/(\lambda_{i}+\eps)$ and dominate the residual diagnostics.
  Assertions and quality checks are made at the recommended operating
  point, not in the $\eps \to 0$ limit.
\item Warm starts: each refinement solve starts from the previous
  iterate, so the inner Krylov iteration counts drop sharply after the
  first solve.
\end{itemize}

\subsection{Relation to augmented Lagrangian methods}

If the gauge were expressed as an explicit constraint $B x = 0$ and the
penalty as $Q = B^{T} W B$, then \eqref{eq:iteration} is algebraically
the Hestenes--Powell augmented-Lagrangian (multiplier) iteration with
fixed penalty parameter \cite{Hestenes1969,Powell1969,FortinGlowinski}:
the update $x^{(k+1)} - x^{(k)} = A_\eps^{-1}(b - A x^{(k)})$ is what a
multiplier update on $Bx = 0$ induces on the primal variable. The
sliding-interface machinery of \texttt{gauged\_fluid.md}~\S5 uses the
same family in its usual role (driving a constraint violation to zero).
The gauge-fixing use is the mirror image: there is no constraint we
want enforced --- any gauge value is acceptable --- and the iteration's
job is to remove the penalty's contamination of the \emph{physical}
modes. The algebra is identical; only the interpretation of which modes
matter is exchanged.

\section{The displacement discontinuity: two spaces, a nodal pairing,
         and a penalised normal jump}
\label{sec:sliding}

The continuous-space gauge of \texttt{gauged\_fluid.md}~\S2 needs no
discontinuity at all for barotropic fluids: tangential interface
continuity is an admissible gauge condition, and one globally continuous
displacement space suffices. A \emph{genuine} tangential jump is needed
when the slip is physical --- two-parameter fluids, and eventually the
non-linear slip map. This section gives the implementation in the detail
the markdown note compresses.

\subsection{The function-space picture}

Let $\Sigma$ be the fluid--solid interface. The discrete field is sought
in the \emph{broken} space
\begin{equation}
  V_{s} \times V_{f}, \qquad
  V_{s} \subset H^{1}(\Omega_{s})^{d}, \quad
  V_{f} \subset H^{1}(\Omega_{f})^{d},
\end{equation}
two entirely independent $H^{1}$ spaces on the two subdomains --- the
discontinuity costs nothing to ``implement'' because nothing couples the
spaces yet. The physics then requires exactly one interface condition on
the displacement, continuity of the normal component,
\begin{equation}
  [\![\, u \cdot n \,]\!] := (u_{s} - u_{f})\cdot n = 0
  \quad \text{on } \Sigma ,
  \label{eq:normaljump}
\end{equation}
with the tangential jump free (the slip). This is a \emph{constraint},
enforced weakly; the tangential freedom needs no action at all. Note the
contrast with a DG discretisation: there is no interior-penalty
stabilisation, no mesh-skeleton terms, and no consistency issue ---
inside each subdomain the space is conforming, and the only broken
surface is $\Sigma$ itself.

\subsection{The nodal pairing $J$}

Enforcing \eqref{eq:normaljump} weakly needs the two interface traces in
one bilinear form, i.e.\ a dictionary between fluid-side and solid-side
interface dofs. This is where MFEM's SubMesh machinery earns its keep.
Both $\Omega_{s}$ and $\Omega_{f}$ are SubMeshes of one parent mesh, and
each space carries a signed dof injection into a parent-mesh space,
\begin{equation}
  \Pi_{s} : V_{s} \to V_{\mathrm{parent}}, \qquad
  \Pi_{f} : V_{f} \to V_{\mathrm{parent}},
\end{equation}
sparse Boolean-up-to-sign matrices (the signs carry the edge/face
orientation bookkeeping for vector-valued nodal bases). A dof of $V_{s}$
and a dof of $V_{f}$ describe \emph{the same interface node} exactly
when their parent images coincide, so the product
\begin{equation}
  J \;=\; \Pi_{s}^{T}\, \Pi_{f}
  \qquad (\text{solid dofs} \times \text{fluid dofs, entries } \pm 1)
  \label{eq:pairing}
\end{equation}
is the complete identification: it is zero except on interface dofs,
where it pairs each fluid trace dof with its solid counterpart,
orientation included. Nothing geometric is computed --- no point
matching, no tolerance --- because the identification is inherited from
the parent's connectivity (\texttt{NewSubMeshPairingMatrix}). In
parallel the same formula runs on \emph{true} dofs,
$J = \Pi_{s}^{T}\Pi_{f}$ as a hypre sparse product
(\texttt{NewSubMeshPairingTrueDofMatrix}); an interface node whose
solid-side and fluid-side owners are \emph{different MPI ranks} is
handled by the product's communication with no special casing.

Because the two trace spaces sit on the same parent faces with the same
nodal basis, they are the \emph{same} discrete trace space in different
coordinates, and $J$ is (up to sign) a permutation between them. That is
what lets every interface bilinear form be assembled \textbf{once}, on
one side only.

\subsection{Why the coupling is weak, not nodal}

Given \eqref{eq:pairing} one is tempted to impose
\eqref{eq:normaljump} \emph{strongly}: constrain the normal component of
each paired dof. This fails at the first corner: on a faceted (or
isoparametric) interface the discrete normal is single-valued on element
faces but \emph{ill-defined at the nodes} where faces meet --- exactly
where the dofs sit. Any nodal choice of normal is arbitrary at the
vertices, and the resulting constraint is inconsistent by $O(h)$.
The weak form avoids the issue entirely, because it only ever evaluates
the normal at quadrature points in face interiors, where it is
unambiguous: with $B$ the normal--normal boundary form on the solid
side,
\begin{equation}
  (B)_{ij} = \int_{\Sigma} (\phi_{i}\cdot n)(\phi_{j}\cdot n)\, dS ,
\end{equation}
the penalty on the normal jump of the paired field expands, via
$u_{f}|_{\Sigma} = $ ($J$-image of the fluid dofs), into the
$2\times 2$ block matrix
\begin{equation}
  \theta\, P, \qquad
  P \;=\;
  \begin{pmatrix}
    B & -B J \\
    -J^{T} B & \;J^{T} B J
  \end{pmatrix},
  \qquad
  \theta\, U^{T} P\, U
  \;=\; \theta \int_{\Sigma} \bigl( (u_{s}-u_{f})\cdot n \bigr)^{2} dS .
  \label{eq:penaltyblock}
\end{equation}
$P$ is symmetric positive semi-definite by construction; its kernel
contains every field with continuous normal trace --- and every purely
tangential jump, which is the point.

\subsection{Exact constraint at moderate $\theta$: the multiplier
iteration}

A pure penalty enforces \eqref{eq:normaljump} only as
$\theta \to \infty$, wrecking the conditioning; this is the same dilemma
as \S\ref{sec:onesolve} in its constraint-enforcing guise, and the same
escape applies --- iterate rather than stiffen. The augmented-Lagrangian
update interleaves with the gauge refinement in a single loop:
\begin{equation}
  \bigl( A + \eps Q + \theta P \bigr)\, U^{(k+1)}
     = f \;-\; w^{(k)} \;+\; \eps\, Q\, U^{(k)} ,
  \qquad
  w^{(k+1)} = w^{(k)} + \theta\, P\, U^{(k+1)} ,
  \label{eq:alupdate}
\end{equation}
starting from $w^{(0)} = 0$. The vector $w$ is the running multiplier:
in the limit it is the discrete \emph{normal traction} on $\Sigma$, the
force the constraint exerts. At a fixed point of \eqref{eq:alupdate} the
multiplier update stalls, which forces $P U = 0$ --- the normal jump
vanishes \emph{exactly} (in the discrete weak sense) at finite
$\theta$ --- and substituting back shows $U$ solves the physical system
with neither $\theta$ nor $\eps$ contaminating it. The standard AL
contraction argument gives a per-sweep reduction of the constraint
violation by a factor $\sim \norm{A_{\mathrm{S}}}/(\norm{A_{\mathrm{S}}}
+ \theta)$ on the relevant Schur complement: moderate $\theta$ and two
or three sweeps replace a stiff penalty, on the same amortised-solver
economics as \S3 (one operator, one preconditioner, several cheap
solves). In practice the normal jump contracts to a floor set by the
interleaved gauge source --- observed at $\sim 6\times 10^{-5}$ of the
normal trace on the coarse test disc, $\theta$-independent --- while
the tangential jump remains $O(1)$: the slip the welded space cannot
represent.

\subsection{The enlarged null space, and what remains open}

A frictionless spherical interface transmits no torque: with the
discontinuity active, the shell and the core rotate \emph{independently},
so the rigid null space is larger than the welded problem's and every
such mode is added to the projector. (Translations of either component
alone are \emph{not} null --- gravity couples them.)

Implemented and verified: the pairing (serial and parallel) and the
penalty/AL machinery on the gravity-free cavity
(\texttt{TestSlidingInterface}, \texttt{TestSlidingInterfacePar});
in the barotropic setting the sliding solution reproduces the condensed
reference and the continuous-space gauge on the solid, with vanishing
normal jump and non-vanishing tangential jump. Still open for the
self-gravitating case: the three-block $[u_{s};\, u_{f};\, \phi]$
solver, and the interface gravity terms when the slip is physical (a
stratified fluid transports its boundary values along $\Sigma$), whose
derivation from the referential energy with the linearised slip map
needs checking before implementation. The mortar (Lagrange-multiplier
trace space) route \cite{Wohlmuth2001} remains the cross-check and the
natural non-linear path.

\section{Algorithm and implementation map}

\paragraph{Algorithm (as implemented).}
\begin{enumerate}
\item Assemble the physical operator $A$ and the penalty $Q$; build the
  solver and preconditioner on $A_{\eps} = A + \eps Q$
  (\texttt{RegularizedMatrix()}).
\item Solve $A_{\eps} x^{(1)} = b$.
\item For $k = 1, \dots, K-1$: form $r = b - A x^{(k)}$ with the
  \emph{physical} $A$; solve $A_{\eps}\, \delta = r$ (warm-started);
  set $x^{(k+1)} = x^{(k)} + \delta$.
\end{enumerate}
For the block (self-gravitating and referential) systems the residual
and update run on the full block operator, with the potential unknowns
accumulated alongside the displacement and the load rows zeroed in the
refinement solves, so the potential converges with the displacement
rather than being recomputed.

\paragraph{Code pointers.}
Base machinery in \texttt{quasi\_static\_problem.hpp}:
\texttt{SetGaugedFluid(marker, $\mu_g$, $\eps$, refinements, penalty)}
with \texttt{GaugePenalty::Deviatoric} or \texttt{::Harmonic},
\texttt{GaugeRefine()}, \texttt{GaugeResiduals()} (per-refinement
residual norms, the diagnostic for \S\ref{sec:semiconv}),
\texttt{SetGaugeEpsilon()}. Block overrides in
\texttt{self\_gravitating.cpp} (gauged fluid) and
\texttt{referential\_problem.cpp} (\texttt{SetVacuumExtension}, the
ball-wide mode kept as a findings log and for the nonlinear future).
Tests: \texttt{TestFluidGauge}, \texttt{TestMixedProblemGauged}
(contraction rates, semi-convergence guards),
\texttt{TestReferentialProblem} (the \S\ref{sec:failure} divergence,
characterised). The decisive external check is the gauged column of the
\texttt{prem\_4} benchmark against \texttt{pyslfp}, including the
degree-0 Love numbers that the Dahlen family cannot reproduce.

\begin{thebibliography}{9}

\bibitem{Lardy1975}
L.~J. Lardy,
\emph{A series representation for the generalized inverse of a closed
linear operator},
Atti Accad. Naz. Lincei Rend. Cl. Sci. Fis. Mat. Natur. (8) 58 (1975),
152--157.
(The origin of iterated Tikhonov.)

\bibitem{EnglHankeNeubauer}
H.~W. Engl, M.~Hanke, A.~Neubauer,
\emph{Regularization of Inverse Problems},
Kluwer, 1996. (Iterated Tikhonov and its qualification; semi-convergence
of iterative regularisation.)

\bibitem{Hestenes1969}
M.~R. Hestenes,
\emph{Multiplier and gradient methods},
J. Optim. Theory Appl. 4 (1969), 303--320.

\bibitem{Powell1969}
M.~J.~D. Powell,
\emph{A method for nonlinear constraints in minimization problems},
in R.~Fletcher (ed.), \emph{Optimization}, Academic Press, 1969.

\bibitem{FortinGlowinski}
M.~Fortin, R.~Glowinski,
\emph{Augmented Lagrangian Methods},
North-Holland, 1983. (Penalty and multiplier iterations for constrained
variational problems, including the fixed-parameter form used here.)

\bibitem{Wohlmuth2001}
B.~I. Wohlmuth,
\emph{Discretization Methods and Iterative Solvers Based on Domain
Decomposition},
Springer, 2001. (Mortar trace spaces: the Lagrange-multiplier
alternative for the interface constraint.)

\end{thebibliography}

\end{document}
